\section{Composition and soundness of the bit comparisons}
\label{sec:composition}

We now compose the inner test with the outer game. The main issue is that the
first player's Fourier truncation depends on the shared background. The joint
form of Lemma~\ref{lem:smoothness} is designed to retain that dependence.

\begin{proposition}[A gap for bit comparisons]\label{prop:bit-gap}
For every fixed integer $J\ge20$, effective choices of constants give a
four-test comparison system with the following properties. It is constructed
from the uniform equation list in Theorem~\ref{thm:max3lin}, for a fixed
completeness error $\zeta>0$ chosen by the construction. Every test compares
two bits in folded proof tables, and its contribution to cost is its failure
probability divided by a positive rational budget $b_j$.
In the YES case there are proof tables of cost at most $4$; in the NO case
every collection of proof tables has cost greater than $J$.
Alphabet sizes are bounded by effective constants depending only on $J$,
and all local numerical parameters depend only on $J$.
\end{proposition}

\begin{proof}
Apply Theorem~\ref{thm:inner} with $J$, taking its atom constants $p,\theta$
to be positive rationals by decreasing them if necessary. Set
\[
 M=mn^{m-1},\qquad q=M-1,\qquad M_*=\frac{\sqrt N}{\sigma}.
\]
The counts $t_0\le t$ of the outer game will be chosen below. Give each actual
first question $U$ a folded bit proof $f_U$ on $A_U$, and each actual second
question $V$ a folded bit proof $f_V$ on $A_V$. The same question always uses
the same proof; an occurrence of a hidden variable is never separately
indexed by its unrevealed equation or position.

The first three tests are those of Theorem~\ref{thm:inner}, applied to $f_V$
under the marginal distribution of $V$. For the fourth test sample a full
outer edge $e$, an independent uniform face array $B$ over $D(A_V)$, and the
Gaussian query table $P=P_{c+\sigma g,\eta\lambda}$ on the common code $\C$.
Compare
\begin{equation}\label{eq:fourth-test}
 f_U\bigl(P(\pi_e^*B)\bigr)
 \quad\hbox{and}\quad
 f_V\bigl(P(B)\bigr).
\end{equation}
Here $\pi_e^*$ acts on every entry of $B$. The Gaussian table on $\C$ is
literally the same on both sides, even if distinct codewords have identical
pullbacks at one alphabet. Let $b_4>0$ be a small rational budget, to be fixed
shortly, and define
\[
 \Cost(f)=\sum_{j=1}^4\frac{\Prob[\text{test $j$ fails}]}{b_j}.
\]

Suppose, toward a contradiction in the NO case, that $\Cost(f)\le J$.
Each of the first three failure bounds required by Theorem~\ref{thm:inner}
then holds, including the random index $V$. Draw an outer edge $e$, a uniform
$i\in[m]$ and $x\in X$, and independent $c$. The $q$ rows other than
$B^i_{x_{-i}}$ will serve as the independent uniform background forms allowed
in Lemma~\ref{lem:hint-soundness}. Put them in a fixed order depending only on
the public $i,x$. The two players receive these rows on their own alphabets,
and use $c,i,x$ as additional public randomness, independent of the game.
From their respective information they form the one-row functions
\begin{align}
 H_V(b)&=h_{x,V}\bigl(B[B^i_{x_{-i}}\leftarrow b],c\bigr),
      &&b\in D(A_V),\label{eq:HV}\\
 H_U(a')&=h_{x,U}\bigl((\pi_e^*B)
                 [B^i_{x_{-i}}\leftarrow a'],c\bigr),
      &&a'\in D(A_U).\label{eq:HU}
\end{align}
The notation in \eqref{eq:HU} simply means filling the missing entry of the
received array with a free argument. The first player needs no knowledge of
$e$ or $V$ to define $H_U$. Both functions are bounded in absolute value by
$M_*$.

\paragraph{The two row functions are close.}
For fixed $e,B,c$, apply \eqref{eq:gaussian-derivative} to the difference
between the two bits in \eqref{eq:fourth-test}, using the same $g,\lambda$.
Bessel's inequality for the orthonormal functions $g_x$ gives
\[
 \frac1N\sum_x
 \abs{h_{x,V}(B,c)-h_{x,U}(\pi_e^*B,c)}^2
 \le\frac1{\sigma^2}\E_{g,\lambda}
 \abs{f_V(P(B))-f_U(P(\pi_e^*B))}^2.
\]
Averaging $e,B,c$ bounds the right-hand side by $4Jb_4/\sigma^2$.
Integrating the missing row in \eqref{eq:HV} restores a fully independent
uniform array $B$, regardless of $i,x$. Thus
\begin{equation}\label{eq:row-comparison}
 \E\norm{H_V-H_U\circ\pi_e^*}_{L^2(D(A_V))}^2
 \le\frac{4Jb_4}{\sigma^2}.
\end{equation}
Choose a positive rational $b_4$ so small that
\begin{equation}\label{eq:budget-choice}
 \frac{64Jb_4}{\sigma^2\theta^2}<\frac p4.
\end{equation}
Markov's inequality then makes the norm in \eqref{eq:row-comparison} less
than $\theta/4$, except on an event of probability less than $p/4$.

\paragraph{Diffuse first-side frequencies remain small after pullback.}
Choose a positive rational $u<1$. At each first-player information state,
let $r_U$ be $H_U$ with all Fourier terms of magnitude at least $u$ removed.
Parseval and the bound $\abs{H_U}\le M_*$ give
\begin{equation}\label{eq:remainder-bounds}
 \#\{\alpha:\abs{\widehat H_U(\alpha)}\ge u\}
 \le\frac{M_*^2}{u^2},\qquad
 \sum_\alpha\abs{\widehat r_U(\alpha)}^4\le u^2M_*^2,
 \qquad
 \norm{r_U}_\infty\le M_r:=M_*+\frac{M_*^2}{u}.
\end{equation}
For the last bound, each removed coefficient satisfies
$\abs{a}\le\abs{a}^2/u$.

We claim that if the total-variation bound of Lemma~\ref{lem:smoothness} for
$d=q+3$ is at most $\Delta$, then
\begin{equation}\label{eq:restricted-fourth}
 \E\sum_{\beta\in D(A_V)^*}
       \abs{\widehat{r_U\circ\pi_e^*}(\beta)}^4
 \le u^2M_*^2+2\Delta M_r^4.
\end{equation}
To prove this, fix $U$ and the independent public randomness $c,i,x$.
For fixed background and $e$, identity \eqref{eq:fourth-moment} expresses the
Fourier sum on the left as the average of
\[
 r_U(\pi_e^*a)r_U(\pi_e^*b)r_U(\pi_e^*c')
 r_U\bigl(\pi_e^*(a+b+c')\bigr),
\]
where $a,b,c'$ are independent uniform forms on $A_V$. Pullback is linear
on the spaces of affine forms, so its fourth argument is determined by the
first three. The entire expression, including the choice of $r_U$, depends
only on $U,c,i,x$ and the $q+3$ pulled-back forms consisting of the background
and these three additional arguments. It is bounded by $M_r^4$ in absolute
value. Replace their joint distribution by independent uniform forms on
$A_U$. Lemma~\ref{lem:smoothness} changes the expectation by at most
$2\Delta M_r^4$. Under the replacement, conditioning on the background fixes
$r_U$, and \eqref{eq:fourth-moment} and \eqref{eq:remainder-bounds} bound its
expectation by $u^2M_*^2$. Averaging proves
\eqref{eq:restricted-fourth}. This argument uses joint total variation; it
does not conditionally replace the missing row while keeping a nonuniform
background fixed.

First choose $u>0$ rational sufficiently small, and then require a sufficiently
small positive rational $\Delta$, so that
\begin{equation}\label{eq:remainder-choice}
 \frac{256\bigl(u^2M_*^2+2\Delta M_r^4\bigr)}{\theta^4}<\frac p4.
\end{equation}
By \eqref{eq:restricted-fourth} and Markov's inequality, no Fourier
coefficient of $r_U\circ\pi_e^*$ has magnitude at least $\theta/4$, except
on an event of probability less than $p/4$.

\paragraph{Local lists give a strategy for the outer game.}
Theorem~\ref{thm:inner} supplies an odd coefficient $\beta$ of $H_V$ with
magnitude at least $\theta$ with probability at least $p$. Combining that
event with the two preceding estimates leaves probability at least $p/2$.
On this event, the coefficient of $H_U\circ\pi_e^*$ at $\beta$ has magnitude
greater than $3\theta/4$, since every coefficient is bounded by the $L^2$
norm of a function. The coefficient of its remainder has magnitude less
than $\theta/4$. Consequently the pullback of the removed part has a nonzero
coefficient at $\beta$. At least one removed frequency $\alpha$ therefore
satisfies
\[
 (\pi_e^*)^*\alpha=\beta.
\]
This conclusion permits multiple frequencies to have the same image:
a nonzero sum has at least one nonzero summand. Since $\pi_e^*(1)=1$ and
$\beta$ is odd, such an $\alpha$ is odd as well.

The first player lists the odd frequencies of $H_U$ of magnitude at least
$u$, and the second lists the odd frequencies of $H_V$ of magnitude at least
$\theta$. These lists are functions of their own information. By
Lemma~\ref{lem:odd-labels}, they are lists of labels, and on the event just
described they contain a pair agreeing under $\pi_e$. Their sizes are at
most $M_*^2/u^2$ and $M_*^2/\theta^2$. Choose uniformly from nonempty lists
and choose any fixed label when a list is empty. With independent private
randomness, the resulting strategy succeeds with probability at least
\begin{equation}\label{eq:decoder-success}
 s_*:=\frac{p/2}{(1+M_*^2/u^2)(1+M_*^2/\theta^2)}>0.
\end{equation}
The strategy is used only to bound the value of a finite game. Its Fourier
computations need not be efficient, and its independent continuous public
randomness does not raise the game value above the maximum deterministic
value.

\paragraph{Parameter order and contradiction.}
All quantities so far depend only on the inner constants and $J$.
Choose $t_0$ effectively so that the upper bound in
Lemma~\ref{lem:hint-soundness} is strictly below $s_*$ for the fixed $q$.
Then choose $t\ge t_0$ so large that Lemma~\ref{lem:smoothness} gives the
required $\Delta$ for $d=q+3$. Equation \eqref{eq:decoder-success}
contradicts the hinted-game upper bound, proving that no NO-case proof has
cost at most $J$.

Finally choose a positive rational $\zeta$ with $t\zeta\le b_4$ and invoke
Theorem~\ref{thm:max3lin}. Its NO bound is always $3/4$, independently of
$\zeta$, so this last choice does not affect the repetition constants.
In the YES case let a global assignment satisfy at least $1-\zeta$ of the
equations. Give each $V$ its values at all requested slots. Give $U$ these
values whenever all its equations are satisfied, and choose any label in
$A_U$ otherwise. The first labels lie in their required affine alphabets,
and the labels agree under $\pi_e$ except on an event of probability at most
$t\zeta$. Use evaluation proofs at these labels. Lemma~\ref{lem:honest}
bounds each of the first three normalized costs by one; agreement of labels
makes the fourth comparison exact, so its normalized cost is also at most
one. This proves the YES bound.

The order of choices is, explicitly,
\[
 J\ \longrightarrow\ \text{inner constants}
 \longrightarrow\ b_4,u,\Delta,s_*
 \longrightarrow\ t_0\ \longrightarrow\ t\ \longrightarrow\zeta.
\]
All strict inequalities can be imposed with rational slack and integer
searches using the displayed bounds. In these searches $M_*^2=N/\sigma^2$
is rational; the possibly irrational quantity $M_r$ may be bounded above
by the rational number $N/\sigma+N/(\sigma^2u)$ before choosing $\Delta$. The complete construction therefore
has effective constants independent of the input length.
\end{proof}
